\begin{proposition}\label{prop:decic-no99}
An elliptic curve over a number field of degree $10$ has no point
of order $99$.
\end{proposition}
\begin{proof}
Put $X=X_1(99)$ and $J=J_1(99)$, of genus $281$.
The complete winding decomposition has weighted dimension $281$.
Every nontrivial-character component has full winding submodule;
the only deficient component has trivial character. Hence the
winding rank-zero criterion gives
\[
                 (\langle2\rangle-1)J(\mathbb Q)\quad\text{finite}.
\]
We do not require $J(\mathbb Q)$ to be finite.
Throughout the proof, diamonds and Hecke operators use point levels:
$T_\ell(E,P)$ is the sum of $(E/C,P\bmod C)$ over the cyclic
subgroups $C$ of order $\ell$.

Let $D$ be the effective rational trace divisor of a hypothetical
level structure over a field of degree $10$. Thus $\deg D=10$,
with multiplicities allowed. The complete finite-field census gives
no noncuspidal places on $X_{\mathbb F_2}$ in degrees at most $9$,
and $120$ in degree $10$. The latter form four diamond orbits.
It follows that $\overline D$ is either entirely cuspidal or a
single noncuspidal degree-$10$ place.

Suppose first that $\overline D$ is entirely cuspidal.
All cusps are defined over $\mathbb Q(\zeta_{99})$, unramified at
$2$, and the cuspidal scheme is etale there. Lift $\overline D$ to
a cuspidal divisor $C$ over a finite unramified extension of
$\mathbb Q_2$. Choose a rational cusp $c_0$ and put
$t=\langle2\rangle$. The class
\[
 (t-1)[D-C]=(t-1)[D-10c_0]+(t-1)[10c_0-C]
\]
is torsion, by the finite rational image and Manin--Drinfeld.
It reduces to zero. The unramified dyadic reduction lemma
\cite[Section~6.2]{DN} shows that this class is killed by $2$.
Consequently
\[
                     2tD+2C\sim2D+2tC.
\]
Both sides have degree $40$, whereas \cite[(2.1)]{NV} gives
\[
 \operatorname{gon}_{\mathbb C}X\geq
 \frac{325\cdot4320}{32768}=\frac{43875}{1024}>40.
\]
The two effective divisors are therefore equal. Separating their
cuspidal and noncuspidal parts gives $tD=D$.
Every point in the support of $D$ has $j$ with a pole at $2$,
since its reduction is cuspidal. In particular its $j$-invariant
is neither $0$ nor $1728$. Its projective diamond stabilizer is
therefore trivial. The projective order of $t$ is $15$, so its
orbits on the support of $D$ all have size $15$. This contradicts
$\deg D=10$.

We are left with the four noncuspidal diamond orbits. We first use
\[
 \pi_{33}:X\longrightarrow X_1(33),\qquad(E,P)\longmapsto(E,3P).
\]
The target has genus $21$ and rank-zero Jacobian by
\cite[Theorem~3.1]{DEHMZ}. Write $\delta_a$ for its point diamond
and put $A_5=T_5-1-5\delta_5$. By Eichler--Shimura and injectivity
of rational torsion reduction at $5$, $A_5$ annihilates its rational
Jacobian. A rational Tate cusp $c$ satisfies $A_5c=0$ as a divisor.
For $P_{33}=\pi_{33*}\overline D$, the degree-zero divisor
\[
                 B_{33}=T_5P_{33}-P_{33}-5\delta_5P_{33}
\]
must therefore be principal. The four exact Riemann--Roch
computations give
\[
 \begin{array}{c|rrrr}
 \text{source orbit}&1&2&3&4\\ \hline
 h^0(B_{33})&0&0&1&0.
 \end{array}
\]
Thus only the third orbit remains. These computations use the six
actual geometric $5$-kernels. Their Frobenius orbit lengths are
$1,1,2,2$, their image places have degrees $10,10,20,20$, and every
residue norm multiplicity is $1$. The Hecke divisor has degree $60$.
A separate native computation and independent marked Velu maps
verify all four tests and their moduli identifications.

For the remaining orbit we use the normalized fibre product
\[
 Y=X_1(9)\mathbin{\times}_{X(1)}X_0(11),\qquad
 \pi(E,P)=(E,11P,\langle9P\rangle).
\]
The curve $Y$ has genus $29$, and $\deg\pi=10$.
The map commutes with Hecke operators and diamonds.
Since $\pi_*\pi^*=[10]$, the finite rational image on $J$ implies
that $(\delta_2-1)J_Y(\mathbb Q)$ is finite. We make no rank-zero
assertion for $J_Y$.
If $E_Y=\pi_*D$, then $[(\delta_2-1)E_Y]$ is rational torsion:
subtract a rational cuspidal divisor of degree $10$ and use
Manin--Drinfeld for the remaining cuspidal difference.
The operator $A_5$ annihilates rational torsion on $J_Y$, again
by reduction at $5$. Thus, for $P_Y=\pi_*\overline D$ and
$T=T_5P_Y$, the divisor
\[
 \begin{split}
 B&=(\delta_2-1)(T_5-1-5\delta_5)P_Y\\
  &=\delta_2T-T+P_Y-\delta_2P_Y
                     +5\delta_5P_Y-5\delta_{10}P_Y
 \end{split}
\]
must be principal. All these operators commute with diamonds,
so one representative of the third source orbit suffices.

We describe the exact computation on $Y$.
On $X_1(9)\simeq\mathbb P^1_t$ take
\[
 c=t^2(t-1),\qquad b=c(t^2-t+1),\qquad
 E_t:y^2+(1-c)xy-by=x^3-bx^2,
\]
with the point $(0,0)$ of order $9$.
Combining its $j$-map with the $j$-map of $X_0(11)$ gives an
irreducible degree-$12$ function field over $\mathbb F_2(t)$.
The saved equations satisfy both original modular equations and
have genus $29$ and exact constant field $\mathbb F_2$.
The functions $t_9$ and $j_{11}$ record the marked $9$-level and
the $j$-invariant of the quotient by the marked cyclic
$11$-subgroup, respectively.

On this model the four $5$-kernel orbits again have lengths
$1,1,2,2$, and all norm multiplicities are $1$.
Both $T$ and $\delta_2T$ have degree $60$.
Since $\delta_{10}P_Y=P_Y$, the positive and negative parts of
$B$ have degree $110$. The exact Riemann--Roch computation gives
\[
                              h^0(B)=0.
\]
A fresh native replay reconstructs the model, verifies all $23$
prime references by exact two-generator ideal identities, and
recomputes the same divisor and Riemann--Roch space.
Independently constructed $5$- and $11$-isogenies identify every
source and Hecke place by the pair $(t_9,j_{11})$, including both
sides of each diamond difference and all residue norm weights.
Hence $B$ is not principal, a contradiction.
\end{proof}
